\section{Experimental Protocol}\label{sec:protocol}

\subsection{Common settings}
Both algorithms are trained under identical conditions:
\begin{itemize}
  \item \textbf{Training budget:} 300,000 environment steps per model, the same for both algorithms; every run trains
        for the full budget (no early stopping).
  \item \textbf{Environment:} the configuration of \cref{sec:env-config} at medium density, identical for both
        algorithms.
  \item \textbf{Network:} an MLP with two hidden layers of 256 units and ReLU for the \ddqn{} Q-network and the \ppo{}
        actor ($25 \to 256 \to 256 \to 5$, 73,733 parameters); the \ppo{} critic is a separate network of the same size.
        $\gam = 0.99$ for both.
  \item \textbf{Seeds:} five training seeds (0--4) per algorithm and reward, so 2 algorithms $\times$ 4 rewards
        $\times$ 5 seeds $=$ 40 models in the main grid. Validation and test episodes use their own disjoint seed sets
        (20 validation seeds, 100 test seeds), and the seeds of the vectorized training environments are derived so
        that they never coincide with them.
  \item \textbf{Compute:} \ddqn{} and \ppo{} were trained on two different machines, one algorithm per machine, with
        the same library versions and at most eight simulator processes running at the same time on each.
\end{itemize}

\subsection{Evaluation protocol}\label{sec:eval-protocol}
\begin{itemize}
  \item \textbf{Validation during training:} every 12,000 environment steps, the current policy is run on the same 20
        validation episodes, giving 25 points per learning curve. The $x$-axis of every curve is environment steps.
  \item \textbf{Final evaluation:} the final checkpoint of each model is evaluated once on 100 test episodes per test
        condition, with the same episode seeds for every model, so all models face identical starting situations.
        Nothing is selected on the test episodes.
  \item \textbf{Policy at evaluation:} deterministic for both algorithms (greedy action for \ddqn{}, most probable
        action for \ppo{}), without learning or exploration.
  \item \textbf{Test conditions:} all evaluation-only (\cref{tab:test-conditions}).
\end{itemize}

\begin{table}[H]
  \centering
  \small
  \caption{Test conditions.}
  \label{tab:test-conditions}
  \begin{tabularx}{\linewidth}{@{}lXX@{}}
    \toprule
    Condition & Differs from training & Purpose \\
    \midrule
    \emph{main}    & nothing (new episodes) & headline result \\
    \emph{low}     & density 0.4 & generalization to easier traffic \\
    \emph{high}    & density 1.0 & generalization to harder traffic \\
    \emph{long}    & 60\,s episodes, 30 vehicles & does the policy keep driving safely for longer? \\
    \emph{precise} & \env{highway-v0}, 15\,Hz physics, collisions between other vehicles
                   & dependence on the coarse physics of the fast variant \\
    \bottomrule
  \end{tabularx}
\end{table}

All models are trained under the same condition (medium density), so \ddqn{} and \ppo{} are compared under identical
training conditions, and their generalization is tested on the unseen low and high densities. Training separate models
at every density would triple the number of runs and is outside our compute budget. Because crash risk per km
grows over a 60\,s episode even for scripted policies, results on \emph{long} are compared with the baselines on the
same condition, not with a model's own \emph{main} result.

\subsection{Metrics}\label{sec:metrics}
All metrics are recorded per episode:
\begin{itemize}
  \item \textbf{Return:} undiscounted episode return under the training reward, and under every other reward
        recomputed from the same episode. Returns are compared only within one reward.
  \item \textbf{Mean speed} (m/s) and \textbf{distance travelled} (m).
  \item \textbf{Completion rate:} share of episodes that reach the time limit without a collision.
  \item \textbf{Collision rate:} share of episodes that end in a collision (at most one collision per episode, since a
        collision terminates the episode); the \textbf{collision-free rate} is its complement. \textbf{Collisions per
        km} are aggregated as total collisions over total distance.
  \item \textbf{Near-collision:} an episode without a collision in which the time-to-collision with the vehicle ahead
        in the ego lane fell below 1.5\,s; we also count the steps with time-to-collision below 1.5\,s. The threshold
        was fixed before training. Only the leader in the ego lane is considered, so cut-ins and conflicts with
        vehicles behind are not captured.
  \item \textbf{Lane behavior:} lane changes per episode (actual changes of the lane index), lane-change reversals
        (a change undone within 3 steps), share of time in the rightmost lane, and the histogram of chosen actions.
  \item \textbf{Overtakes per km:} vehicles ahead of the ego at the start and behind it at the end, divided by the
        distance.
\end{itemize}

\subsubsection{Hardware-independent computational cost}
Wall-clock times cannot be compared between the algorithms, because \ddqn{} and \ppo{} were trained on different
machines. We therefore compare quantities that do not depend on the hardware:
\begin{itemize}
  \item \textbf{Environment steps to a target performance} and \textbf{to the plateau} of the learning curve. Two
        target levels were fixed before training, each requiring a safety and a speed condition on the validation
        episodes at three consecutive validation points: \emph{beats\_idle} (collision rate below 0.6 and mean speed
        above 24.8\,m/s, i.e.\ strictly better than idle on both) and \emph{strong} (collision rate at most 0.4 and
        mean speed at least 25.5\,m/s). No scripted baseline reaches either level. The plateau is computed on the
        validation return under the run's own reward: the final level is the mean over the last 20\% of the
        validation points, the improvement is the final level minus the first validation value, and the plateau point
        is the earliest validation point from which the rolling mean over five points stays within $0.1 \times$ the
        improvement of the final level. A plateau point inside the last 20\% of the curve is reported as censored.
  \item \textbf{Gradient computation:} the number of gradient updates and of samples processed by them per
        environment step.
  \item \textbf{Inference cost:} at evaluation both policies are a single forward pass through the same
        $25 \to 256 \to 256 \to 5$ network (the \ddqn{} Q-network or the \ppo{} actor), i.e.\ 73,216
        multiply-accumulate operations per decision, so their inference cost is identical by construction.
\end{itemize}
Wall-clock training time and measured inference time (median over 1,000 single-observation calls on one CPU thread)
are reported per machine, together with the machine description, and are used only to compare runs of the same
algorithm.

\subsection{Hyper-parameter selection}\label{sec:hp-selection}
Shared settings (budget, network, $\gam$, validation schedule) are fixed for both algorithms. Each algorithm's own
hyper-parameters were tuned with the same budget: at most four configurations (the defaults plus up to three
variations), one run each with a dedicated tuning seed that is never used for results, on the balanced reward at
medium density, evaluated on the validation episodes only. The configuration with the highest mean balanced return
over the last three validation points was selected, ties broken by the lower collision rate. Hyper-parameter
sensitivity is then studied one parameter at a time around the selected values, with three seeds per setting on the
balanced reward: the learning rate at 1/3 and 3 times its value for both algorithms, the duration of the
$\varepsilon$ decay $\pm 0.2$ of the budget for \ddqn{}, and the clipping parameter $\pm 0.1$ for \ppo{}. The base
setting is the corresponding main-grid run, so the sensitivity study adds 24 runs. Tuning used a single seed and only the balanced reward, and the selected values are used for all
rewards; this is a limitation of the budget. Final values for both algorithms are in \cref{app:hyperparameters}.

\subsection{Statistical methodology}\label{sec:stats}
\paragraph{State of the art.}
Results of deep RL depend strongly on random seeds and implementation details~\cite{henderson2018matters}.
Reporting only the mean over a few runs can be misleading; robust aggregate statistics such as the interquartile
mean with bootstrapped confidence intervals are recommended~\cite{agarwal2021rliable}.

\paragraph{Our protocol.}
The training seed is the unit of analysis. For each model we average the metrics over its 100 test episodes; across
the five seeds we report the mean with a two-sided 95\% Student-$t$ interval ($t = 2.776$ for five seeds) and always
show the individual runs next to it. We use the $t$-interval rather than a bootstrap because a bootstrap over five
values cannot leave the observed range and underestimates the uncertainty when the standard deviation itself is
estimated from so few runs. Baselines have no training seed, so their intervals are Wilson intervals over episodes.
Five seeds support only coarse conclusions~\cite{patterson2024empirical}; differences within the intervals are not
interpreted.

\subsection{Experiments}
\Cref{tab:experiments} maps the experiments to the subsections of \cref{sec:results}.

\begin{table}[H]
  \centering
  \small
  \caption{Experiments.}
  \label{tab:experiments}
  \begin{tabularx}{\linewidth}{@{}XXl@{}}
    \toprule
    Experiment & Runs & Results \\
    \midrule
    \ddqn{} vs.\ \ppo{}, learning curves, efficiency, safety
      & main grid: 2 algorithms $\times$ 4 rewards $\times$ 5 seeds $=$ 40
      & \cref{sec:res-ddqn-ppo,sec:res-curves,sec:res-efficiency,sec:res-safety} \\
    Reward design & the same 40 models, compared across rewards & \cref{sec:res-reward} \\
    Traffic density & the same 40 models on low and high density & \cref{sec:res-density} \\
    Hyper-parameter sensitivity & 24 runs (4 settings $\times$ 3 seeds per algorithm) & \cref{sec:res-hp} \\
    Qualitative behavior & videos of the same test episodes for both algorithms, early and final checkpoints
      & \cref{sec:res-qualitative} \\
    Computational cost & all runs & \cref{sec:res-compute} \\
    \bottomrule
  \end{tabularx}
\end{table}
